\documentclass{article}
\usepackage{mathtools} 
\usepackage{fontspec}
\usepackage[UTF8]{ctex}
\usepackage{amsthm}
\usepackage{mdframed}
\usepackage{xcolor}
\usepackage{amssymb}
\usepackage{amsmath}
\usepackage{tikz}
\usepackage{pgfplots}
\usepgfplotslibrary{polar}
\usetikzlibrary{3d} % 允许3D坐标
% \pgfplotsset{compat=1.18}


% 定义新的带灰色背景的说明环境 zremark
\newmdtheoremenv[
  backgroundcolor=gray!10,
  % 边框与背景一致，边框线会消失
  linecolor=gray!10
]{zremark}{说明}

% 通用矩阵命令: \flexmatrix{矩阵名}{元素符号}{行数}{列数}
\newcommand{\flexmatrix}[4]{
  \[
  #1 = \begin{pmatrix}
    #2_{11}     & #2_{12}     & \cdots & #2_{1#4}   \\
    #2_{21}     & #2_{22}     & \cdots & #2_{2#4}   \\
    \vdots      & \vdots      & \ddots & \vdots     \\
    #2_{#31}    & #2_{#32}    & \cdots & #2_{#3#4}
  \end{pmatrix}
  \]
}

% 简化版命令（默认矩阵名为A，元素符号为a）: \quickmatrix{行数}{列数}
\newcommand{\quickmatrix}[2]{\flexmatrix{A}{a}{#1}{#2}}

\begin{document}
\title{习题 11.2}
\author{张志聪}
\maketitle

\section*{1}

\begin{itemize}
  \item 定理11.2（比较原则）

        因为$g(x)$是$[a, + \infty)$上的非负函数，
        于是对任意$u_1, u_2 \in [a, + \infty)$，且$u_1 < u_2$，有
        \begin{align*}
          \int_{a}^{u_1} g(x) dx - \int_{a}^{u_2} g(x) dx
           & = \int_{u_1}^{u_2} g(x) dx
        \end{align*}
        因为$x \in [u_1, u_2]$上，$g(x) \geq 0$，于是由积分不等式可知
        \begin{align*}
          \int_{u_1}^{u_2} g(x) dx \geq 0
        \end{align*}
        所以，$G(u) = \int_{a}^{u} g(x) dx$是单调递增函数。
        同理可证，$F(u) = \int_{a}^{u} f(x) dx$也是单调递增函数。

        已知$\int_{a}^{+\infty} g(x) dx = \lim \limits_{u \to +\infty} \int_{a}^{u} g(x) dx$收敛，
        即$\lim \limits_{u \to +\infty} G(u)$收敛，设收敛于$J$，$G(u)$单调递增，所以$J$是$G(u)$的上确界。

        因为任意有限区间$[a, u]$上，都有$f(x) \leq g(x)$，由积分不等式可知
        \begin{align*}
          \int_{a}^{u} f(x) dx \leq \int_{a}^{u} g(x) dx
        \end{align*}
        即
        \begin{align*}
          F(u) \leq G(u) \leq J
        \end{align*}
        所以，$J$也是$F(u)$的上界。由单调有界定理可知，$\lim \limits_{u \to +\infty} F(u)$收敛，
        即$\int_{a}^{+\infty} f(x) dx = \lim\limits_{u \to +\infty} \int_{a}^{u} f(x) dx = \lim \limits_{u \to +\infty} F(u)$收敛。

  \item 推论1

        \begin{itemize}
          \item (i)

                $\lim\limits_{x \to +\infty} \frac{f(x)}{g(x)} = c$，
                由极限定义可知，对$\epsilon = \frac{1}{2} c$，存在$M > 0$，当$x > M$，有
                \begin{align*}
                  \left|\frac{f(x)}{g(x)} - c \right| < \epsilon = \frac{1}{2} c \\
                  \frac{1}{2}c < \frac{f(x)}{g(x)} < \frac{3}{2} c               \\
                  \frac{c}{2} g(x) < f(x) < \frac{3c}{2} g(x)
                \end{align*}
                (1) $\frac{c}{2} g(x) < f(x)$如果$\int_{a}^{+\infty} f(x) dx$收敛，
                由比较原则，$\int_{a}^{+\infty} \frac{c}{2} g(x) dx$收敛，
                又
                \begin{align*}
                  \int_{a}^{+\infty} \frac{c}{2} g(x) dx
                   & = \lim\limits_{u \to +\infty} \int_{a}^{u} \frac{c}{2} g(x) dx \\
                   & = \lim\limits_{u \to +\infty} \frac{c}{2} \int_{a}^{u} g(x) dx \\
                   & = \frac{c}{2} \lim\limits_{u \to +\infty} \int_{a}^{u} g(x) dx
                \end{align*}
                所以$\lim\limits_{u \to +\infty} \int_{a}^{u} g(x) dx $收敛，即
                $\int_{a}^{+\infty} g(x) dx$收敛。

                (2)$f(x) < \frac{3c}{2} g(x)$如果$\int_{a}^{+\infty} f(x) dx$发散，
                由比较原则，$\int_{a}^{+\infty} \frac{3c}{2} g(x) dx$发散，
                从而$\int_{a}^{+\infty} g(x) dx$发散。


                同理可证，由于$f(x) < \frac{3c}{2} g(x)$则$g(x)$收敛，$f(x)$收敛；
                由于$\frac{c}{2} g(x) < f(x)$则$g(x)$发散，$f(x)$发散。

          \item (ii)(iii) 略
        \end{itemize}
\end{itemize}

\section*{2}

 (i)

由$\frac{a + b}{2} \geq \sqrt{ab} (a, b \geq 0)$，有
\begin{align*}
  \frac{1}{2}f^2(x) + g^2(x) \geq |f(x) g(x)|
\end{align*}
已知$\int_{a}^{+\infty} f(x) dx, \int_{a}^{+\infty} g(x) dx$收敛，
由比较原则可得$\int_{a}^{+\infty} |f(x) g(x)| dx$收敛，
绝对收敛可得
\begin{align*}
  \int_{a}^{+\infty} f(x) g(x) dx
\end{align*}
收敛.

(ii)

\begin{align*}
  [f(x) + g(x)]^2 = f^2(x) + g^2(x) + 2 f(x) g(x)
\end{align*}
已知$\int_{a}^{+\infty} f^2(x) dx, \int_{a}^{+\infty} g^2(x) dx, \int_{a}^{+\infty} 2 f(x) g(x) dx$收敛，
于是由极限的四则运算法则可知
\begin{align*}
  \int_{a}^{+\infty} [f(x) + g(x)]^2 dx
   & = \int_{a}^{+\infty} f^2(x) + g^2(x) + 2 f(x) g(x) dx
\end{align*}
收敛.

\section*{3}

\begin{itemize}
  \item (1)

        已知$h(x) \leq f(x) \leq g(x)$，所以
        \begin{align*}
          0 \leq g(x) - f(x) \leq g(x) - h(x)
        \end{align*}
        $\because$ $\int_{a}^{+\infty} h(x) dx, \int_{a}^{+\infty} g(x) dx$收敛.\\
        $\therefore$ $\int_{a}^{+\infty} g(x) - h(x) dx$收敛.

        由比较原则可得
        \begin{align*}
          \int_{a}^{+\infty} g(x) - f(x) dx
        \end{align*}
        收敛。我们有
        \begin{align*}
          \int_{a}^{+\infty} f(x) dx & =  - \int_{a}^{+\infty} [g(x) - f(x)] - g(x) dx
        \end{align*}
        由$\int_{a}^{+\infty} g(x) - f(x) dx$和$\int_{a}^{+\infty} g(x) dx$收敛，可得
        $\int_{a}^{+\infty} f(x) dx $收敛。

  \item (2)

        令
        \begin{align*}
          H(u) & = \int_{a}^{u} h(x) dx \\
          F(u) & = \int_{a}^{u} f(x) dx \\
          G(u) & = \int_{a}^{u} g(x) dx
        \end{align*}
        已知$h(x) \leq f(x) \leq g(x)$，所以由积分不等式
        \begin{align*}
          \int_{a}^{u} h(x) dx \leq \int_{a}^{u} f(x) dx \leq \int_{a}^{u} g(x) dx
        \end{align*}
        即
        \begin{align*}
          H(u) \leq F(u) \leq G(u)
        \end{align*}
        已知$\int_{a}^{+\infty} h(x) dx, \int_{a}^{+\infty} g(x) dx$收敛，
        即$\lim \limits_{u \to +\infty} H(u) = \lim \limits_{u \to +\infty} G(u) = A$，
        由迫敛性可知
        \begin{align*}
          \lim \limits_{u \to +\infty} F(u) = A
        \end{align*}
        即
        \begin{align*}
          \int_{a}^{+\infty} f(x) dx = A
        \end{align*}
\end{itemize}

\section*{4}

\begin{itemize}
  \item (1)
        \begin{align*}
          \lim\limits_{x \to +\infty} x^{\frac{4}{3}} \frac{1}{\sqrt[3]{x^4 + 1}} = 1
        \end{align*}
        由柯西判别法（推论3）可知
        \begin{align*}
          \int_{a}^{+\infty} \frac{1}{\sqrt[3]{x^4 + 1}} dx
        \end{align*}
        收敛。

  \item (2)

        考虑
        \begin{align*}
          - \frac{x}{1 - e^x} = \frac{x}{e^{x} - 1}
        \end{align*}
        于是$\frac{x}{e^{x} - 1}$是$[1, +\infty)$上的非负函数。

        我们有
        \begin{align*}
          \lim\limits_{x \to +\infty} x^2 \cdot \frac{x}{e^{x} - 1} = 0
        \end{align*}
        （洛必达法则，容易得到结果）

        由柯西判别法（推论3）可知
        \begin{align*}
          \int_{1}^{+\infty} \frac{x}{e^{x} - 1} dx
        \end{align*}
        收敛，所以
        \begin{align*}
          - \int_{1}^{+\infty} \frac{x}{e^{x} - 1} dx
           & = \int_{1}^{+\infty} \frac{x}{1 - e^x} dx
        \end{align*}
        收敛。

  \item (3)

        提示:
        \begin{align*}
          \lim \limits_{x \to +\infty} x^{\frac{1}{2}} \cdot \frac{1}{1 + \sqrt{x}} = 1
        \end{align*}
        发散。

  \item (4)

        提示
        \begin{align*}
          \lim \limits_{x \to +\infty} x^2 \cdot \frac{x arctan\, x}{1 + x^3} = \frac{\pi}{2}
        \end{align*}
        收敛。

  \item (5)

        $n > 1$，取$p = \frac{1 + p}{2} > 1$，于是$n - p > 1$
        \begin{align*}
          \lim\limits_{x \to +\infty} x^{p} \cdot \frac{\ln (1 + x)}{x^n}
           & = \lim\limits_{x \to +\infty} \frac{\ln (1 + x)}{x^{n - p}} \\
           & = 0
        \end{align*}
        所以无穷积分收敛。

        $n = 1$，有
        \begin{align*}
          \lim\limits_{x \to +\infty} x \cdot \frac{\ln (1 + x)}{x}
           & = \lim\limits_{x \to +\infty} \ln (1 + x) \\
           & = +\infty
        \end{align*}
        所以无穷积分发散。

        $n < 1$，有
        \begin{align*}
          \frac{\ln (1 + x)}{x^n} > \frac{\ln (1 + x)}{x}
        \end{align*}
        由比较原则，无穷积分发散。

        综上，$n \leq 1$，无穷积分发散；
        $n > 1$，无穷积分收敛。

  \item (6)

        $n - m > 1$，有
        \begin{align*}
          \lim \limits_{x \to +\infty} x^{n - m} \cdot \frac{x^m}{1 + x^n} = 1
        \end{align*}
        由柯西判别法（推论3）可知，无穷积分收敛。

        $n - m = 1$，有
        \begin{align*}
          \lim \limits_{x \to +\infty} x^{n - m} \cdot \frac{x^m}{1 + x^n} = 1
        \end{align*}
        由柯西判别法（推论3）可知，无穷积分发散。

        $n - m < 1$，有
        \begin{align*}
          \frac{x^m}{1 + x^n} > \frac{x^m}{1 + x^{m + 1}}
        \end{align*}
        由比较原则可知，无穷积分发散。

        综上，$n - m > 1$，无穷积分收敛；
        $n - m \leq 1$，无穷积分发散。
\end{itemize}

\section*{5}

\begin{itemize}
  \item (1)

        令$t = \sqrt{x}, dx = 2tdt$，$x: 1 \to +\infty, t: 1 \to +\infty$，于是
        \begin{align*}
          \int_{1}^{+\infty} \frac{sin \sqrt{x}}{x} dx
           & = \int_{1}^{+\infty} 2t \frac{sin t}{t^2} dt \\
           & = 2 \int_{1}^{+\infty} \frac{sin t}{t} dt
        \end{align*}
        由例3可知，无穷积分条件收敛。

  \item (2)
        \begin{align*}
          \left| \frac{sgn(sin\, x)}{1 + x^2} dx \right|
           & \leq \frac{1}{1 + x^2}
        \end{align*}
        又因为
        \begin{align*}
          \lim\limits_{x \to +\infty} x^2 \cdot \frac{1}{1 + x^2} = 1
        \end{align*}
        由柯西判别法（推论3）可知，无穷积分$\int_{0}^{+\infty} \frac{1}{1 + x^2} dx$收敛。
        由比较原则可知
        \begin{align*}
          \int_{0}^{+\infty} \left| \frac{sgn(sin\, x)}{1 + x^2} dx \right| dx
        \end{align*}
        收敛。

        综上，原无穷积分绝对收敛。

  \item (3)

        (i) 条件收敛。
        \begin{align*}
          \frac{\sqrt{x} cos\, x}{100 + x}
           & = cos\, x \cdot \frac{\sqrt{x}}{100 + x}
        \end{align*}
        因为
        \begin{align*}
          \int_{0}^{+\infty} cos\, x dx
        \end{align*}
        有界。
        \begin{align*}
          \lim\limits_{x \to +\infty} \frac{\sqrt{x}}{100 + x} = 0
        \end{align*}
        且
        \begin{align*}
          (\frac{\sqrt{x}}{100 + x})'
           & = \frac{\frac{1}{2\sqrt{x}}(100 + x) - \sqrt{x}}{(100 + x)^2} \\
           & = \frac{100 - x}{2\sqrt{x}(100 + x)^2}                        \\
        \end{align*}
        于是$x \geq 0$时，单调递增。

        由定理11.3（狄利克雷判别法）可知，$\int_{100}^{+\infty} \frac{\sqrt{x} cos\, x}{100 + x} dx$
        收敛。又
        \begin{align*}
          \int_{0}^{+\infty} \frac{\sqrt{x} cos\, x}{100 + x} dx
           & = \int_{0}^{100} \frac{\sqrt{x} cos\, x}{100 + x} dx + \int_{100}^{+\infty} \frac{\sqrt{x} cos\, x}{100 + x} dx
        \end{align*}
        且$\int_{0}^{100} \frac{\sqrt{x} cos\, x}{100 + x} dx$为初等函数的定积分，所以
        $\int_{0}^{+\infty} \frac{\sqrt{x} cos\, x}{100 + x} dx$收敛。

        (ii) 非绝对收敛。

        \begin{align*}
          \left|\frac{\sqrt{x} cos\, x}{100 + x} \right|
           & \geq \frac{\sqrt{x} cos^2\, x}{100 + x}                                              \\
           & = \frac{\frac{1}{2} \sqrt{x} (1 + cos2\, x) }{100 + x}                               \\
           & = \frac{\frac{1}{2} \sqrt{x} }{100 + x} + \frac{\frac{1}{2} \sqrt{x} cos2x}{100 + x}
        \end{align*}
        因为
        \begin{align*}
          \lim \limits_{x \to +\infty}\sqrt{x} \cdot \frac{\frac{1}{2} \sqrt{x} }{100 + x} = \frac{1}{2}
        \end{align*}
        于是$\int_{100}^{+\infty} \frac{\frac{1}{2} \sqrt{x} }{100 + x} dx$发散。
        且利用狄利克雷判别法可知，$\int_{100}^{+\infty} \frac{\frac{1}{2} \sqrt{x} cos2x}{100 + x} dx$收敛。
        于是
        \begin{align*}
          \int_{100}^{+\infty} \frac{\sqrt{x} cos\, x}{100 + x} dx
           & = \int_{100}^{+\infty} \frac{\frac{1}{2} \sqrt{x} }{100 + x} dx
          + \int_{100}^{+\infty} \frac{\frac{1}{2} \sqrt{x} cos2x}{100 + x} dx
        \end{align*}
        发散。

        综上
        \begin{align*}
          \int_{0}^{+\infty} \frac{\sqrt{x} cos\, x}{100 + x} dx
           & = \int_{0}^{100} \frac{\sqrt{x} cos\, x}{100 + x} dx + \int_{100}^{+\infty} \frac{\sqrt{x} cos\, x}{100 + x} dx
        \end{align*}
        由于右侧第一个积分是初等函数的定积分。故$\int_{0}^{+\infty} \frac{\sqrt{x} cos\, x}{100 + x} dx $发散。


  \item (4)

        (i) 条件收敛

        \begin{align*}
          \int_{e}^{+\infty} sin x dx
        \end{align*}
        有界。
        \begin{align*}
          \lim\limits_{x \to +\infty} \frac{\ln(\ln x)}{\ln x}
           & = \lim\limits_{x \to +\infty} \frac{\frac{1}{\ln x}\frac{1}{x}}{\frac{1}{x}} \\
           & = \lim\limits_{x \to +\infty} \frac{1}{\ln x}                                \\
           & = 0
        \end{align*}
        又
        \begin{align*}
          \left[\frac{\ln(\ln x)}{\ln x}\right]'
           & = \frac{\frac{1}{\ln x}\frac{1}{x} \ln x - \ln(\ln x) \frac{1}{x}}{\ln^2 x} \\
           & = \frac{\frac{1}{x}(1 - \ln(\ln x))}{\ln^2 x}
        \end{align*}
        所以$x \geq e^e$时，单调递减，
        于是由定理11.3（狄利克雷判别法）可得
        \begin{align*}
          \int_{e^e}^{+\infty} \frac{\ln(\ln x)}{\ln x} sin x dx
        \end{align*}
        收敛。我们有
        \begin{align*}
          \int_{e}^{+\infty} \frac{\ln(\ln x)}{\ln x} sin x dx
           & = \int_{e}^{e^e} \frac{\ln(\ln x)}{\ln x} sin x dx + \int_{e^e}^{+\infty} \frac{\ln(\ln x)}{\ln x} sin x dx
        \end{align*}
        右侧第一个是定积分，所以$\int_{e}^{+\infty} \frac{\ln(\ln x)}{\ln x} sin x dx$收敛。

        (ii) 非绝对收敛。

        \begin{align*}
          \left|\frac{\ln(\ln x)}{\ln x} sin x \right|
           & \geq \frac{\ln(\ln x)}{\ln x} sin^2 x                                               \\
           & = \frac{1}{2} \frac{\ln(\ln x)}{\ln x} (1 - cos2x)                                  \\
           & = \frac{1}{2} \frac{\ln(\ln x)}{\ln x} - \frac{1}{2} \frac{\ln(\ln x)}{\ln x} cos2x
        \end{align*}

        因为
        \begin{align*}
          \lim \limits_{x \to +\infty} x \cdot \frac{\ln(\ln x)}{\ln x}
           & = \lim \limits_{x \to +\infty} \frac{\ln(\ln x) + x \frac{1}{\ln x}\frac{1}{x}}{\frac{1}{x}} \\
           & = \lim \limits_{x \to +\infty} x \ln(\ln x) + \frac{x}{\ln x}                                \\
           & = +\infty
        \end{align*}
        由柯西判别法可知，$\int_{e^e}^{+\infty} \frac{\ln(\ln x)}{\ln x} dx$ 发散。

        利用狄里克莱判别法，可知$\int_{e^e}^{+\infty} \frac{\ln(\ln x)}{\ln x} cos2x dx$收敛。
        所以$\int_{e^e}^{+\infty} \frac{\ln(\ln x)}{\ln x} sin^2 x dx$发散。我们有
        \begin{align*}
          \int_{0}^{+\infty} \frac{\ln(\ln x)}{\ln x} sin^2 x dx
           & = \int_{0}^{e^e} \frac{\ln(\ln x)}{\ln x} sin^2 x dx + \int_{e^e}^{+\infty} \frac{\ln(\ln x)}{\ln x} sin^2 x dx
        \end{align*}
        右侧第一个是定积分，所以$\int_{0}^{+\infty} \frac{\ln(\ln x)}{\ln x} sin^2 x dx$发散。

        综上，该无穷积分非绝对收敛。

\end{itemize}

\section*{6}

\begin{itemize}
  \item (1)

        定义
        \begin{align*}
          \int_{1}^{+\infty} \frac{sin\, x}{\sqrt{x}} dx \\
          \int_{1}^{+\infty} \frac{sin^2\, x}{x} dx
        \end{align*}
        第一个无穷积分，借助例3的结论。

  \item (2)

        定义
        \begin{equation*}
          f(x) =
          \begin{cases}
            0,   & [n - 1, n -  \frac{1}{2^n \cdot 2^n}) \\
            2^n, & [n - \frac{1}{2^n \cdot 2^n}, n)
          \end{cases}
          n = 1, 2, 3, ...
        \end{equation*}
        于是
        \begin{align*}
          \int_{1}^{+\infty} f(x) dx
           & = \lim\limits_{x \to +\infty} \int_{1}^{u} f(x) dx                                     \\
           & = \lim\limits_{n \to +\infty} \sum\limits_{i = 1}^n 2^i \frac{1}{2^i \cdot 2^i}        \\
           & = \lim\limits_{n \to +\infty} \sum\limits_{i = 1}^n \frac{1}{2^i}                      \\
           & = \lim\limits_{n \to +\infty} \frac{\frac{1}{2}(1 - (\frac{1}{2})^n)}{1 - \frac{1}{2}} \\
           & = \lim\limits_{n \to +\infty} 1 - (\frac{1}{2})^n                                      \\
           & = 1
        \end{align*}
        收敛。

        又
        \begin{align*}
          \int_{1}^{+\infty} f^2(x) dx
           & = \lim\limits_{x \to +\infty} \int_{1}^{u} f^2(x) dx                                      \\
           & = \lim\limits_{n \to +\infty} \sum\limits_{i = 1}^n 2^i \cdot 2^i \frac{1}{2^i \cdot 2^i} \\
           & = \lim\limits_{n \to +\infty} \sum\limits_{i = 1}^n 1                                     \\
           & = \lim\limits_{n \to +\infty} n                                                           \\
           & = + \infty
        \end{align*}
        发散。
\end{itemize}

\section*{7}
\begin{itemize}
  \item 方法一：

        已知$\lim\limits_{x \to +\infty} f(x) = 0$，于是我们有
        \begin{align*}
          c & = \lim\limits_{x \to +\infty} \frac{f^2(x)}{|f(x)|} \\
            & = \lim\limits_{x \to +\infty} |f(x)|                \\
            & = 0
        \end{align*}
        注：以上由极限定义容易验证。

        已知$\int_{a}^{+\infty} f(x) dx$绝对收敛，
        则$\int_{a}^{+\infty} |f(x)| dx$收敛，
        由比较原则的推理1-(ii)可知
        \begin{align*}
          \int_{a}^{+\infty} f^2(x) dx
        \end{align*}
        收敛。

  \item 方法二：

        已知$\lim\limits_{x \to +\infty} f(x) = 0$，则
        \begin{align*}
          \lim\limits_{x \to +\infty} |f(x)| = 0
        \end{align*}
        于是对$0 < \epsilon < 1$，存在$M > 0$，$x > M$，有
        \begin{align*}
          |f(x)| < \epsilon
        \end{align*}
        已知$0 < \epsilon < 1$，此时有
        \begin{align*}
          |f(x)| < f^2(x)
        \end{align*}
        因为$\int_{a}^{+\infty} f(x) dx$绝对收敛，
        由比较原则可知
        \begin{align*}
          \int_{a}^{+\infty} f^2(x) dx
        \end{align*}
        注：$\int_{a}^{M} f^2(x) dx$是定积分，不影响敛散性。

\end{itemize}

\section*{8}


\end{document}